\chapter{Deep Q-Networks: De Tablas a Redes Neuronales}
\label{cap:dqn}

\section{La Maldición de la Dimensionalidad y la Necesidad de Generalizar}
En el Capítulo~\ref{cap:td} consolidamos el algoritmo Q-Learning tabular. Asignar un número real a cada par estado-acción en una matriz bidimensional es intuitivo y convergente, pero resulta inviable en la robótica moderna:
\begin{itemize}
  \item Si un robot percibe una imagen de cámara de $84 \times 84$ píxeles en escala de grises con 256 niveles de brillo, el número de estados posibles es $256^{84 \times 84} \approx 10^{16973}$, una cifra enormemente superior al número de átomos en el universo observable.
  \item Si el robot mide magnitudes físicas continuas (como la velocidad en m/s o el ángulo de giro en grados), el espacio de estados es incontable: la probabilidad de visitar exactamente el mismo estado real dos veces es prácticamente cero.
\end{itemize}

La solución consiste en sustituir la tabla rígida por una \textbf{función paramétrica diferenciable}:
\begin{equation}
  Q(s, a; \boldsymbol{\theta}) \approx Q^*(s, a)
\end{equation}
donde $\boldsymbol{\theta}$ representa el vector de pesos sinápticos de una red neuronal artificial profunda. La ventaja colosal de las redes neuronales es su capacidad de \textbf{generalización inductiva}: la red aprende patrones geométricos y conceptuales, infiriendo valores $Q$ precisos para estados que jamás presenció durante el entrenamiento basándose en su similitud con experiencias anteriores.

\begin{optibot}
OptiBot recibe una actualización completa de hardware. Se eliminan las celdas discretas de la cuadrícula y se equipa con:
\begin{enumerate}
  \item Un sensor LiDAR frontal continuo que mide la distancia al obstáculo más cercano en metros ($d \in [0.05, 20.0]$ m).
  \item Un sensor de batería que reporta el porcentaje continuo exacto de carga ($b \in [0.0\%, 100.0\%]$).
  \item Un encoder de odometría que registra la velocidad lineal actual ($v \in [0.0, 3.0]$ m/s).
  \item Una brújula digital que calcula el ángulo relativo hacia el muelle de descarga ($\theta \in [-\pi, \pi]$ rad).
\end{enumerate}
El estado es un vector continuo cuatridimensional $\mathbf{s} = [d, b, v, \theta]^T \in \mathbb{R}^4$. Las acciones siguen siendo discretas: $\{ \text{Avanzar}, \text{Girar Izquierda}, \text{Girar Derecha}, \text{Recargar} \}$.
\end{optibot}

\section{La Tríada Mortal de Sutton \& Barto}
Durante más de dos décadas, entrenar redes neuronales profundas con algoritmos de aprendizaje por refuerzo fue considerado impracticable debido a divergencias numéricas e inestabilidad extrema. Richard Sutton y Andrew Barto identificaron la causa teórica subyacente, conocida como la \textbf{Tríada Mortal}. La divergencia matemática acecha siempre que coinciden simultáneamente estos tres factores:
\begin{enumerate}
  \item \textbf{Aproximación de funciones no lineal:} El uso de redes neuronales en lugar de tablas discretas exactas.
  \item \textbf{Bootstrapping:} Utilizar la predicción del propio modelo en $t+1$ como objetivo para corregir la predicción en $t$ (la esencia de Bellman y TD).
  \item \textbf{Aprendizaje fuera de política (\textit{Off-Policy}):} Aprender sobre la política óptima codiciosa a partir de datos generados por una política exploratoria diferente.
\end{enumerate}

\section{Los Dos Pilares Estabilizadores de DQN (DeepMind 2015)}
En 2015, Volodymyr Mnih y el equipo de Google DeepMind publicaron en la revista \textit{Nature} el algoritmo \textbf{Deep Q-Network (DQN)}, demostrando por primera vez que una red neuronal profunda podía aprender a jugar decenas de videojuegos de Atari al nivel de un experto humano directamente desde los píxeles crudos. El éxito de DQN radica en dos mecanismos esenciales que desactivan la Tríada Mortal:

\subsection{1. Buffer de Repetición de Experiencias (\textit{Experience Replay})}
En el aprendizaje supervisado estándar, las muestras del conjunto de entrenamiento son independientes e idénticamente distribuidas (i.i.d.). En RL, por el contrario, los datos llegan en secuencias temporales fuertemente correlacionadas: el estado en $t+1$ depende críticamente de lo que el robot hizo en $t$. Entrenar una red neuronal con datos secuenciales provoca \textbf{olvido catastrófico} y oscilaciones destructivas.

El \textbf{Replay Buffer} $\mathcal{D}$ almacena las últimas $N$ transiciones observadas por el robot:
\[
  e_t = (s_t, a_t, r_{t+1}, s_{t+1}, \text{done}_{t+1})
\]
Durante el entrenamiento, en lugar de actualizar los pesos con la transición inmediata, el algoritmo extrae un \textbf{mini-batch aleatorio uniforme} de $\mathcal{D}$. Esto rompe por completo las correlaciones temporales y permite reutilizar cada experiencia valiosa en múltiples actualizaciones de gradiente.

\subsection{2. Red Objetivo Congelada (\textit{Target Network})}
En el Q-Learning tabular, modificar la celda $Q(s,a)$ no alteraba el valor numérico de las demás celdas. Sin embargo, en una red neuronal, un paso de retropropagación que actualiza los pesos $\boldsymbol{\theta}$ para corregir $Q(s, a)$ altera sutilmente las salidas de la red para \textbf{todos los demás estados y acciones}. Dado que el objetivo TD depende de la propia red ($\max_{a'} Q(s', a'; \boldsymbol{\theta})$), actualizar $\boldsymbol{\theta}$ equivale a ``perseguir un blanco móvil que se aleja con cada disparo''.

DQN desacopla el objetivo manteniendo dos redes idénticas:
\begin{itemize}
  \item La \textbf{Red Online} ($\boldsymbol{\theta}$): Se actualiza en cada paso mediante descenso de gradiente.
  \item La \textbf{Red Objetivo o Target Network} ($\boldsymbol{\theta}^-$): Sus pesos permanecen fijos y congelados durante $C$ pasos. Se utiliza exclusivamente para calcular el valor del objetivo TD:
  \begin{equation}
    y_i = r + \gamma (1 - \text{done}) \max_{a'} Q(s', a'; \boldsymbol{\theta}^-)
  \end{equation}
\end{itemize}
Cada $C$ pasos temporales (por ejemplo, cada 1,000 pasos), los pesos de la red online se copian a la red objetivo: $\boldsymbol{\theta}^- \leftarrow \boldsymbol{\theta}$.

\section{Función de Pérdida de Bellman y Double DQN}
La función de pérdida que se minimiza mediante descenso de gradiente es el \textbf{Error Cuadrático Medio de Bellman}:
\begin{equation}
  \mathcal{L}(\boldsymbol{\theta}) = \mathbb{E}_{(s, a, r, s', d) \sim \mathcal{D}} \left[ \Big( r + \gamma (1 - d) \max_{a'} Q(s', a'; \boldsymbol{\theta}^-) - Q(s, a; \boldsymbol{\theta}) \Big)^2 \right]
\end{equation}

\subsection{Double DQN (Hado van Hasselt, 2016)}
El operador $\max_{a'}$ en el objetivo de DQN induce un \textbf{sesgo de sobreestimación} sistemático: si el valor Q de una acción contiene ruido estocástico positivo, el operador de máximo lo seleccionará preferentemente, inflando artificialmente el valor estimado del futuro.

\textbf{Double DQN} elimina este sesgo desacoplando la selección de la acción de su evaluación:
\begin{equation}
  y_i^{\text{DDQN}} = r + \gamma (1 - d) Q\left( s', \, \arg\max_{a'} Q(s', a'; \boldsymbol{\theta}); \, \boldsymbol{\theta}^- \right)
\end{equation}
La red online $\boldsymbol{\theta}$ selecciona la mejor acción codiciosa, pero la red objetivo $\boldsymbol{\theta}^-$ evalúa objetivamente cuánto vale dicha acción.

\section{Implementación en PyTorch: Red Neuronal y Entrenamiento}
\begin{lstlisting}[style=rlcode]
import torch
import torch.nn as nn
import torch.optim as optim
import random
from collections import deque

class QNetwork(nn.Module):
    def __init__(self, state_dim=4, action_dim=4):
        super().__init__()
        self.net = nn.Sequential(
            nn.Linear(state_dim, 64),
            nn.ReLU(),
            nn.Linear(64, 64),
            nn.ReLU(),
            nn.Linear(64, action_dim)
        )
    def forward(self, x):
        return self.net(x)

class ReplayBuffer:
    def __init__(self, capacity=50000):
        self.buffer = deque(maxlen=capacity)
    def push(self, s, a, r, s_next, done):
        self.buffer.append((s, a, r, s_next, done))
    def sample(self, batch_size):
        batch = random.sample(self.buffer, batch_size)
        s, a, r, s_next, d = zip(*batch)
        return (torch.FloatTensor(s), torch.LongTensor(a).unsqueeze(1),
                torch.FloatTensor(r).unsqueeze(1), torch.FloatTensor(s_next),
                torch.FloatTensor(d).unsqueeze(1))
    def __len__(self):
        return len(self.buffer)

# Paso de entrenamiento de DQN con Target Network
def train_step(online_net, target_net, optimizer, buffer, batch_size=32, gamma=0.99):
    if len(buffer) < batch_size:
        return 0.0
    states, actions, rewards, next_states, dones = buffer.sample(batch_size)
    
    with torch.no_grad():
        # Double DQN: seleccion con online_net, evaluacion con target_net
        best_actions = online_net(next_states).argmax(1, keepdim=True)
        next_q = target_net(next_states).gather(1, best_actions)
        targets = rewards + gamma * (1 - dones) * next_q
        
    current_q = online_net(states).gather(1, actions)
    loss = nn.MSELoss()(current_q, targets)
    
    optimizer.zero_grad()
    loss.backward()
    optimizer.step()
    return loss.item()
\end{lstlisting}

\begin{notabox}
\textbf{La culminación técnica:} OptiBot dispone ahora de un cerebro artificial con aprendizaje por refuerzo profundo capaz de procesar telemetría continua y generalizar ante situaciones nunca vistas. En el capítulo final, pondremos a prueba este sistema en un proyecto integrador completo y analizaremos las fronteras que extienden estos cimientos al estado del arte contemporáneo.
\end{notabox}
